\section{Energy and the retarded force}\label{ret:section}

The estimates below use the conserved energy and two different changes
of variables on a backward light cone. We first derive the force kernels
without derivatives of the density. This representation belongs to the
Glassey--Strauss approach~\cite{GlasseyStrauss1986}; the organization by
cone-controlled field components also appears in
\cite[Proposition~2.2 and Section~3]{LukStrain2014Continuation}.
We give the source-label calculation explicitly, including its initial
boundary term.

Throughout this section, fix a classical solution on a compact interval
$[0,t_*]$, with $t_*\leq T_0<\infty$, in the class specified in the main
theorem. All limiting operations are justified on this fixed compact
interval. Constants in the final bounds depend only on the initial
data and $T_0$; constants used solely to remove a cutoff may also depend
on the compact-interval momentum support. Retain the characteristic notation fixed in Section~\ref{sec:setup}.

\begin{lemma}[Transport and cone energy]\label{ret:energy}
The characteristic flow preserves phase volume and transports $f_0$.
In particular $0\leq f\leq\|f_0\|_\infty$. Define
\[
 \mathcal U=\frac{|E|^2+|B|^2}{2}+\int qf\,dv,
 \qquad
 \mathcal S=E\times B+\int quf\,dv,
 \qquad
 \mathcal H_0=\int_{\R^3}\mathcal U(0,x)\,dx.
\]
Then
\begin{equation}\label{eq:energy}
 \int_{\R^3}\mathcal U(s,x)\,dx=\mathcal H_0
 \qquad(0\leq s\leq t_*).
\end{equation}
For every vertex $(t,x)$ in this interval, put
$r=t-s$, $y=x-rn$, $n\in S^2$, and
\[
 d=1-n\cdot u(v),\qquad
 G=|E(s,y)+n\times B(s,y)|+|n\cdot B(s,y)|.
\]
Then
\begin{equation}\label{eq:cone-energy}
 \int_0^t\int_{S^2}r^2
 \left(G^2+\int qd f(s,x-rn,v)\,dv\right)dn\,ds
 \leq4\mathcal H_0.
\end{equation}
\end{lemma}

\begin{proof}
The phase vector field has zero divergence: $u$ is independent of $x$,
$E$ is independent of $v$, and
$\nabla_v\cdot(u\times B)=B\cdot(\nabla_v\times\nabla_v q)=0$.
Liouville's formula gives volume preservation, and the Vlasov equation
gives $f(s,Y,V)=f_0(z)$. The flow exists backwards from every finite
phase point on this compact time interval. Indeed, its spatial path
stays in a bounded ball because $|u|<1$; the fields are bounded on the
resulting compact spacetime set, so its momentum cannot escape in
finite time.

Multiplication of Vlasov by $q$ and integration in $v$, followed by
Maxwell's energy identity, give
\[
 \partial_s\mathcal U+\nabla_x\cdot\mathcal S=0,
 \qquad |\mathcal S|\leq\mathcal U.
\]
There is no momentum boundary term because $f$ has compact phase-space
support on the interval. Integrate against spatial cutoffs equal to one
on a ball of radius $R$, with gradient bounded by $C/R$. The error tends
to zero: the fields are bounded in $C([0,t_*];L^2)$ and the particle
energy has compact support and is continuous in time. Passing to the
limit proves conservation.

The outward spatial normal of the ball $|y-x|<t-s$ is $-n$. Its energy
derivative is therefore minus the boundary integral of
$\mathcal U-n\cdot\mathcal S$. At the cone tip the ball energy tends
to zero by local continuity. Hence the exact cone identity is
\[
 \int_0^t\int_{S^2}r^2
   (\mathcal U-n\cdot\mathcal S)(s,x-rn)\,dn\,ds
 =\int_{|y-x|<t}\mathcal U(0,y)\,dy.
\]
Its particle part is $\int qd f\,dv$, and its field part is
\[
 \frac12\bigl(|E+n\times B|^2+|n\cdot B|^2\bigr).
\]
Since $(a+b)^2\leq2(a^2+b^2)$, the last identity implies
\eqref{eq:cone-energy}.
\end{proof}

\subsection{The retarded source labels}

For fixed $t,x$, set $r=|x-y|$, $s=t-r$, and $n=(x-y)/r$ when $r>0$.
The zero-data wave potentials are
\begin{equation}\label{ret:potentials}
 (\Phi,\mathcal A)(t,x)=\frac1{4\pi}
  \int_{|x-y|<t}\int_{\R^3}
       \frac{f(t-|x-y|,y,v)}{|x-y|}(1,u(v))\,dv\,dy.
\end{equation}

The following change to source labels also appears in
\cite[Section~9, Lemma~9.2]{Lazarovici2016}. We give the calculation
for the present flow.

\begin{lemma}[Retarded phase change]\label{ret:phase-change}
Away from the cone tip, the map sending $(y,v)$ at time
$s=t-|x-y|$ to its initial label is one-to-one, with Jacobian
\begin{equation}\label{ret:phase-jacobian}
 dz=d\,dy\,dv,\qquad d=1-n\cdot u(v)>0.
\end{equation}
Its image consists, up to the null set of tip labels, of
\[
 \mathcal L_{t,x}=\{z:|x-Y(0,z)|<t\}.
\]
For such a label, the retarded time solves
\begin{equation}\label{ret:time-equation}
 s+|x-Y(s,z)|=t.
\end{equation}
Consequently \eqref{ret:potentials} equals
\begin{equation}\label{ret:label-potentials}
 (\Phi,\mathcal A)(t,x)=\frac1{4\pi}
       \int_{\mathcal L_{t,x}} f_0(z)\frac{(1,u(V(s,z)))}{rd}\,dz.
\end{equation}
\end{lemma}

\begin{proof}
Let $Z_s$ be the fixed-time inverse phase flow. Its determinant is one
by Lemma~\ref{ret:energy}, and
$\partial_sZ_s=-D Z_s\,(u,V')$. Since $\nabla_y s=n$, the derivative
of $Z_{t-|x-y|}(y,v)$ is the fixed-time derivative multiplied by
\[
 \Id-(u,V')\otimes(n,0).
\]
The determinant of this rank-one perturbation is $1-n\cdot u=d$.

Along a fixed label, the function on the left of
\eqref{ret:time-equation} is strictly increasing. More explicitly,
on any compact part of that characteristic, its increment between
$s_1<s_2$ is at least
$(1-\sup|u|)(s_2-s_1)>0$. It starts below $t$ exactly for labels in
$\mathcal L_{t,x}$ and at $s=t$ is at least $t$. Thus it has exactly
one root. The root is a tip only if $Y(t,z)=x$. This set has phase
Lebesgue measure zero, by the volume-preserving fixed-time flow and
Fubini's theorem. This proves injectivity and the image assertion.
Transport of $f$ and \eqref{ret:phase-jacobian} now give
\eqref{ret:label-potentials}.
\end{proof}

\subsection{The exact electric and magnetic kernels}

Along a source characteristic write
\begin{equation}\label{ret:source-acceleration}
 b=u'=\frac{\Id-u\otimes u}{q}(E+u\times B),
\end{equation}
where these fields, $u$, and $q$ are evaluated at its retarded point.

\begin{proposition}[Retarded field representation]\label{ret:representation}
Let $E^{\rm fr},B^{\rm fr}$ be the homogeneous wave solutions with data
\[
 (E^{\rm fr},\partial_tE^{\rm fr})|_{t=0}
       =(E_0,\nabla\times B_0),\qquad
 (B^{\rm fr},\partial_tB^{\rm fr})|_{t=0}
       =(B_0,-\nabla\times E_0).
\]
For $t>0$, define initial boundary terms by
\begin{equation}\label{ret:initial-boundary}
 \begin{split}
 (E^{\rm bd},B^{\rm bd})(t,x)
 =\frac1{4\pi t}\int_{|x-y_0|=t}\int_{\R^3}
 \frac{f_0(y_0,v_0)}{d_0}
 \bigl(n_0-u_0,\ n_0\times(n_0-u_0)\bigr)\,dv_0\,dS_{y_0},\\
 n_0=(x-y_0)/t,\qquad u_0=u(v_0),\qquad d_0=1-n_0\cdot u_0.
 \end{split}
\end{equation}
Set these terms to zero at $t=0$. Then
\begin{equation}\label{ret:fields}
 \begin{split}
 E(t,x)&=E^{\rm fr}(t,x)+E^{\rm bd}(t,x)
                  +\frac1{4\pi}\int_{\mathcal L_{t,x}}f_0(z)\,\mathscr K\,dz,\\
 B(t,x)&=B^{\rm fr}(t,x)+B^{\rm bd}(t,x)
                  +\frac1{4\pi}\int_{\mathcal L_{t,x}}f_0(z)\,(n\times\mathscr K)\,dz,
 \end{split}
\end{equation}
where
\begin{equation}\label{ret:kernel}
 \mathscr K=\frac{n-u}{r^2q^2d^3}
       +\frac{(n-u)(n\cdot b)-db}{rd^3}.
\end{equation}
All these integrals are well defined at the stated classical regularity.
Moreover,
\begin{equation}\label{ret:data-bound}
 \sup_{0\leq t\leq T_0,\ x\in\R^3}
 (|E^{\rm fr}|+|B^{\rm fr}|+|E^{\rm bd}|+|B^{\rm bd}|)
 \leq C_{T_0,f_0,E_0,B_0}.
\end{equation}
\end{proposition}

\begin{proof}
\emph{Wave representation.}
Kirchhoff's formula gives
$\Box(\Phi,\mathcal A)=(\rho_f,j_f)$ with zero data, where
$\Box=\partial_t^2-\Delta$. Maxwell's equations give, distributionally,
\[
 \Box E=-\nabla\rho_f-\partial_tj_f,\qquad
 \Box B=\nabla\times j_f.
\]
Thus the fields are
$(-\nabla\Phi-\partial_t\mathcal A,\nabla\times\mathcal A)$ plus
the stated homogeneous waves. Here it is important that the electric
potential contribution has initial time derivative $-j_f(0)$.

We justify this assertion with only $C^1$ sources. In the variables
$y=x-rn$, the potentials integrate a $C^1$ source against $r\,dr\,dn$.
Their first derivatives are continuous, and locally uniformly in $x$,
\[
 \partial_t\mathcal A(t,x)=t j_f(0,x)+O(t^2),\qquad
 \nabla\Phi(t,x),\ \nabla\mathcal A(t,x)=O(t^2).
\]
These formulas give the required zero initial values and initial
difference quotients of the potential fields. For completeness,
uniqueness of the resulting continuous distributional wave remainder
can be reduced to classical uniqueness: convolve spatially with a
compact smooth mollifier. The wave equation then supplies a continuous
second time derivative, including at zero; the initial difference
quotients identify its zero initial velocity. Local wave energy gives
zero, and removal of the mollifier gives the assertion. No second
derivative of $f$ is used.

\emph{Interior kernels.}
To differentiate \eqref{ret:label-potentials}, hold the initial label
fixed. Implicit differentiation of \eqref{ret:time-equation} gives
\[
 s_t=1/d,\qquad\nabla_xs=-n/d.
\]
Write $\ell=(rd)^{-1}$. With $s$ held fixed,
\[
 \nabla_x\ell=-\frac{n-u}{r^2d^2}.
\]
With $x$ fixed,
\[
 r_s=-n\cdot u,\qquad
 n_s=-\frac{u-n(n\cdot u)}r,\qquad
 \ell_s=\frac{q^{-2}-d}{r^2d^2}+\frac{n\cdot b}{rd^2}.
\]
Therefore the interior electric kernel is
\[
 -\nabla_x^{\rm total}\ell-\partial_t(u\ell)
 =\frac{n-u}{r^2d^2}+\frac{n-u}{d}\ell_s-\frac b{rd^2},
\]
which simplifies to \eqref{ret:kernel}. In the magnetic curl the
corresponding expression is its cross product with $n$: indeed
\[
 \nabla_x^{\rm total}\times(u\ell)
 =\left(-\frac{n-u}{r^2d^2}-\frac n d \ell_s\right)\times u
                  -\frac{n\times b}{rd^2}
 =n\times\mathscr K.
\]

\emph{Boundary terms and the cone tip.}
The label domain has indicator $\mathbf1_{\{t>|x-y_0|\}}$.
Its time derivative is $\delta(t-|x-y_0|)$ and its spatial gradient is
$-n_0\delta(t-|x-y_0|)$. The resulting electric factor is
$(n_0-u_0)/(td_0)$ and the magnetic one is its cross product with $n_0$.
Integration over the initial sphere gives \eqref{ret:initial-boundary}.
The initial velocities are bounded on $\operatorname{supp}f_0$, so
$d_0$ has a fixed positive lower bound. Sphere area and compact initial
velocity support consequently give
$|E^{\rm bd}|+|B^{\rm bd}|\leq C_{f_0}t$, uniformly in $x$.

The differentiations at the cone tip can be justified by inserting a
smooth cutoff that vanishes for $r<\gamma$ and equals one for
$r>2\gamma$. On the supported labels of the fixed compact time interval,
$d$ is bounded below and $u,b$ are bounded. Derivatives of this cutoff
are at most $C/\gamma$. By \eqref{ret:phase-jacobian}, the omitted
potential and cutoff errors are bounded by integrals of $Cf/r$ over
$r<2\gamma$, multiplied by at most $C/\gamma$ for the latter. They
are $O(\gamma^2)$ and $O(\gamma)$, respectively. The omitted kernel
integrals are $O(\gamma)$ because their worst radial singularity after
this change of variables is $r^{-2}$. These bounds are locally uniform
for $t>0$. The already established initial traces handle $t=0$.
Thus no term supported at the cone tip remains.

Finally, Kirchhoff's formula for a homogeneous wave $W$ with data
$(W_0,W_1)$ bounds it by
\[
 \|W(t)\|_\infty
 \leq\|W_0\|_\infty+t\|\nabla W_0\|_\infty+t\|W_1\|_\infty.
\]
The assumed bounded initial derivatives give \eqref{ret:data-bound}.
The homogeneous waves here need not together be a vacuum Maxwell
solution; the representation requires only their displayed wave data.
\end{proof}

\subsection{Receiver geometry and pointwise force bounds}

Let $(X(t),V_X(t))$ be a supported receiver characteristic and write
\[
 \begin{gathered}
 a=u(V_X),\qquad q_X=q(V_X),\qquad K=V_X',\\
 D=1-n\cdot a,\qquad e=1-a\cdot u,\qquad Q=D\Id+n\otimes a.
 \end{gathered}
\]
Since $h+a\times(n\times h)=Qh$, Proposition~\ref{ret:representation}
and \eqref{ret:phase-jacobian} give
\begin{equation}\label{ret:force}
 K=K^{\rm data}+\frac1{4\pi}\int_{|X(t)-y|<t}\int_{\R^3}
                (T_{\rm ker}+S_{\rm ker})\,dv\,dy,
\end{equation}
where $|K^{\rm data}|\leq C_{T_0,f_0,E_0,B_0}$ and
\begin{equation}\label{ret:force-kernels}
 T_{\rm ker}=\frac{fQ(n-u)}{r^2q^2d^2},\qquad
 S_{\rm ker}=\frac{fQ\bigl((n-u)(n\cdot b)-db\bigr)}{rd^2}.
\end{equation}
The density and source quantities in these integrals are evaluated at
$(t-r,y,v)$, not at the receiver.

Figure~\ref{fig:retarded-geometry} shows the source and receiver geometry.
\begin{figure}[tbp]
  \centering
  \begin{tikzpicture}[
      x=0.96cm,y=0.96cm,font=\small,
      >={Stealth[length=2mm]},
      source/.style={draw=blue!65!black,line width=1pt},
      receiver/.style={draw=orange!55!black,line width=1pt,densely dashed},
      event/.style={circle,fill=black,inner sep=1.7pt},
      guide/.style={draw=black!28,densely dotted},
      labelbox/.style={fill=white,inner sep=1.5pt}
    ]
    \node at (3.4,5.2) {\textbf{(a)} Spacetime: schematic $1+1$ slice};
    \fill[blue!4] (0.55,0.9)--(3.55,3.9)--(6.55,0.9)--cycle;
    \draw[black!38,dashed] (0.55,0.9)--(3.55,3.9)--(6.55,0.9);
    \draw[->,black!65] (0.15,0.4)--(0.15,4.75) node[above] {time};
    \draw[->,black!65] (0.15,0.4)--(6.7,0.4) node[right] {space};
    \draw[guide] (0.15,2.05)--(1.7,2.05);
    \draw[guide] (0.15,3.9)--(3.55,3.9);
    \node[left] at (0.15,2.05) {$s$};
    \node[left] at (0.15,3.9) {$t$};
    \draw[source] (1.2,0.45)
      .. controls (1.3,1.1) and (1.525,1.7) .. (1.7,2.05)
      .. controls (2.0,2.65) and (2.12,3.6) .. (2.35,4.5);
    \draw[receiver] (2.8,0.45)
      .. controls (2.86,1.45) and (2.9,2.2) .. (3.16,2.95)
      .. controls (3.316,3.4) and (3.525,3.8) .. (3.55,3.9)
      .. controls (3.6075,4.13) and (3.65,4.3) .. (3.72,4.5);
    \node[source,draw=none,anchor=east] at (2.17,4.42) {source $Y$};
    \node[receiver,draw=none,anchor=west] at (3.86,4.42) {receiver $X$};
    \draw[->,line width=1.1pt] (1.7,2.05)--(3.55,3.9);
    \node[event] at (1.7,2.05) {};
    \node[event] at (3.55,3.9) {};
    \node[anchor=south east,inner sep=0pt] at (1.60,2.23) {$(s,Y(s))$};
    \node[labelbox,anchor=west] at (3.7,3.87) {$(t,X(t))$};
    \node[align=center,black!60] at (4.85,1.42) {backward\\light cone};
    \node[labelbox,rotate=45] at (2.58,3.18) {light ray};

    \node at (10.85,5.2) {\textbf{(b)} Spatial geometry};
    \node[black!65] at (10.85,4.68) {Positions at times $s$ and $t$};
    \draw[->,line width=1.1pt] (8.3,2.4)--(12.65,2.4)
      node[midway,above=4pt] {$n$};
    \draw[source,->] (8.3,2.4)--(9.3,3.72)
      node[above=3pt] {$u(V(s))$};
    \draw[receiver,->] (12.65,2.4)--(13.36,3.72)
      node[above=3pt] {$a(t)$};
    \node[event] at (8.3,2.4) {};
    \node[event] at (12.65,2.4) {};
    \node[below=5pt] at (8.3,2.4) {$Y(s)$};
    \node[below=5pt] at (12.65,2.4) {$X(t)$};
    \draw[|<->|,black!65] (8.3,1.45)--(12.65,1.45)
      node[midway,fill=white,inner sep=3pt] {$r$};
    \node at (10.7,0.66) {$\displaystyle
      r=t-s,\qquad n=\frac{X(t)-Y(s)}{r}$};
  \end{tikzpicture}
  \caption{Retarded geometry along a receiver characteristic.
  The source event $(s,Y(s))$ lies on the backward light cone of
  $(t,X(t))$, with $r=t-s=|X(t)-Y(s)|>0$.
  The spatial unit vector $n$ points from $Y(s)$ to $X(t)$.
  Panel~(b) shows the source velocity $u(V(s))=Y'(s)$ and the
  receiver velocity $a(t)=X'(t)$ at these respective events.
  The curves and vector directions are schematic.}
  \label{fig:retarded-geometry}
\end{figure}


\begin{lemma}[Pointwise force bounds]\label{ret:direct-bounds}
Let $\theta,\phi>0$ satisfy $\theta^2\asymp d$ and $\phi^2\asymp D$,
with fixed comparison constants. The transport kernel splits into two
terms, according to the normal and tangential parts of $n-u$, bounded by
\begin{equation}\label{ret:transport-bounds}
 T_r=\frac{Cf}{r^2q^2\theta^2},\qquad
 T_\perp=\frac{Cf\phi}{r^2q^2\theta^3}.
\end{equation}
With $G$ as in Lemma~\ref{ret:energy}, the source acceleration satisfies
\begin{equation}\label{ret:acceleration-bounds}
 |b|\leq\frac Cq(G+\theta^2|B|),\qquad
 |n\cdot b|\leq\frac{C\theta}q(G+\theta^2|B|),\qquad
 |q'|\leq C(G+\theta|B|).
\end{equation}
Consequently
\begin{equation}\label{ret:source-bounds}
 |S_{\rm ker}|\leq S_g+S_b,\qquad
 S_g=\frac{Cf\phi}{rq\theta^2}G,\qquad
 S_b=\frac{Cf\phi}{rq}|B|.
\end{equation}
If $V_X\ne0$ and $P_X$ is orthogonal projection onto $V_X^\perp$, then
$|P_XS_{\rm ker}|\leq C\phi(S_g+S_b)$.
\end{lemma}

\begin{proof}
The exact identities
\[
 2d=|n-u|^2+q^{-2},\qquad
 2D=|n-a|^2+q_X^{-2}
\]
give $|n-u|\leq C\theta$ and $|n-a|\leq C\phi$.
Also $Qn=n$, whereas for $h\perp n$,
\[
 |Qh|\leq D|h|+|a\cdot h|\leq C\phi|h|.
\]
The normal component of $n-u$ is $dn$ and its tangential component has
norm at most $C\theta$. This proves \eqref{ret:transport-bounds}.

Put $u=(1-d)n+u_\perp$ and $H=E+n\times B$. The force at the source is
\[
 F_{\rm src}=E+u\times B=H+(-dn+u_\perp)\times B.
\]
Its tangential and normal components satisfy
\[
 |(F_{\rm src})_\perp|\leq C(G+\theta^2|B|),\qquad
 |n\cdot F_{\rm src}|\leq C(G+\theta|B|).
\]
Indeed the tangential part of $u_\perp\times B$ involves only
$n\cdot B$, while its normal part is bounded by $C\theta|B|$.
For $L_u=\Id-u\otimes u$ we have
\[
 |L_un|\leq C\theta,\qquad
 |n\cdot L_un|=2d-d^2\leq C\theta^2,
 \qquad |n\cdot L_uh|\leq C\theta|h|\quad(h\perp n).
\]
Combining these component bounds with $b=q^{-1}L_uF_{\rm src}$ proves
the first two inequalities in \eqref{ret:acceleration-bounds}.
The last follows from $q'=u\cdot E=u\cdot H-u\cdot(n\times B)$.

The numerator $N=(n-u)(n\cdot b)-db$ is tangent to $n$, since
$n\cdot N=d(n\cdot b)-d(n\cdot b)=0$. This cone tangency is also used
in \cite[Lemma~2.1]{Pallard2015Refined}. The preceding estimates give
\[
 |N|\leq\frac{C\theta^2}{q}(G+\theta^2|B|).
\]
The tangential bound for $Q$ now proves \eqref{ret:source-bounds}.
Finally $P_Xa=0$, so $|P_Xn|=|P_X(n-a)|\leq C\phi$. For $h\perp n$,
\[
 |P_XQh|\leq D|h|+|P_Xn|\,|a\cdot h|
               \leq C\phi^2|h|.
\]
Applying this to $N$ gives the additional projected factor.
\end{proof}

\begin{lemma}[Source time along a receiver]\label{ret:receiver-time}
Along a retarded source label satisfying
$t=s+|X(t)-Y(s,z)|$ away from collisions,
\begin{equation}\label{ret:receiver-jacobians}
 \frac{dt}{ds}=\frac dD,\qquad
 ds\,f_0(z)\,dz=D\,dt\,f(t-r,y,v)\,dy\,dv.
\end{equation}
For $f_0(z)dz$-almost every label there is no same-time collision
$Y(t,z)=X(t)$ anywhere in the compact time interval. Such a label has
one continuously differentiable retarded branch after its initial-cone entry, when that
entry occurs, and $r$ is bounded below on each closed part of this
branch in the compact interval.
\end{lemma}

\begin{proof}
Differentiate $t-s-|X(t)-Y(s,z)|=0$ to obtain $D\,dt=d\,ds$.
Both factors are positive. Combining this with
\eqref{ret:phase-jacobian} gives the measure identity in
\eqref{ret:receiver-jacobians}.

To exclude same-time collisions, let $R_v$ bound the supported momenta
on $[0,t_*]$. Use a time grid of mesh at most $\zeta$ with
$O(1+t_*/\zeta)$ points. If a collision occurs, at a nearest grid time
the two positions are at distance at most $2\zeta$, since both spatial
speeds are at most one. At each grid time the mass of labels satisfying
this distance bound is at most
$C\|f_0\|_\infty R_v^3\zeta^3$, by phase-volume preservation.
The union has mass $O(\zeta^2)$ and tends to zero.

For the remaining labels, strict timelikeness gives uniqueness and
continuous differentiability of the retarded root. The entry time is unique because
$t-|X(t)-Y(0,z)|$ is strictly increasing. At entry $r=t>0$, since
an entry at zero would be an excluded initial collision. On any
closed part of the branch, continuity and the absence of collisions
give a positive minimum of $r$. This assertion is labelwise; no
uniform lower bound over all labels is used.
\end{proof}
